% ============================================================
% DOCUMENTCLASS
% ============================================================
\documentclass[12pt,oneside]{book}

% ============================================================
% METADATOS
% ============================================================
\newcommand{\universidad}{Universidad de Buenos Aires (UBA)}
\newcommand{\materia}{Derecho Civil --- Parte General: Personas Jurídicas}
\newcommand{\titulo}{Personas Jurídicas}
\newcommand{\autor}{Isaac Edgar Camacho Ocampo}
\newcommand{\anio}{2024 -- 2025}

% ============================================================
% PAQUETES
% ============================================================
\usepackage[utf8]{inputenc}
\usepackage[T1]{fontenc}
\usepackage[spanish,es-nodecimaldot]{babel}

\usepackage[
    top=2.5cm,
    bottom=2.5cm,
    left=3cm,
    right=2.5cm,
    headheight=15pt
]{geometry}

\usepackage{amsmath}
\usepackage{amssymb}
\usepackage{mathtools}

\usepackage{graphicx}
\usepackage{float}
\usepackage{caption}
\usepackage{subcaption}

\usepackage{booktabs}
\usepackage{array}
\usepackage{longtable}
\usepackage{tabularx}

\usepackage{enumitem}

\usepackage{xcolor}
\usepackage[most]{tcolorbox}

\usepackage{fancyhdr}

\usepackage{ragged2e}

\graphicspath{
    {./img/}
    {./graficos/}
    {./figuras/}
}

\usepackage[
    colorlinks=true,
    linkcolor=blue!50!black,
    citecolor=green!40!black,
    urlcolor=blue!60!black,
    pdftitle={\titulo},
    pdfauthor={\autor},
    pdfsubject={Apuntes universitarios},
    bookmarks=true
]{hyperref}

% ============================================================
% CONFIGURACIÓN
% ============================================================
\setcounter{tocdepth}{2}

\setlength{\parindent}{0pt}
\setlength{\parskip}{0.5em}

\setlist[itemize]{
    topsep=0.3em,
    itemsep=0.2em
}

\setlist[enumerate]{
    topsep=0.3em,
    itemsep=0.2em
}

\clubpenalty=10000
\widowpenalty=10000

\pagestyle{fancy}
\fancyhf{}

\fancyhead[L]{\nouppercase{\leftmark}}
\fancyhead[R]{\materia}

\fancyfoot[C]{\thepage}

\renewcommand{\headrulewidth}{0.4pt}
\renewcommand{\footrulewidth}{0pt}

\fancypagestyle{plain}{
    \fancyhf{}
    \fancyfoot[C]{\thepage}
    \renewcommand{\headrulewidth}{0pt}
}

% ============================================================
% COMANDOS
% ============================================================
\newcommand{\abs}[1]{\left|#1\right|}
\newcommand{\norm}[1]{\left\lVert#1\right\rVert}
\newcommand{\vect}[1]{\mathbf{#1}}
\newcommand{\deriv}[2]{\frac{d#1}{d#2}}
\newcommand{\pderiv}[2]{\frac{\partial#1}{\partial#2}}

% ============================================================
% ENTORNOS / CAJAS
% ============================================================
\newtcolorbox{definition}[1]{
    enhanced,
    breakable,
    title={Definición: #1},
    fonttitle=\bfseries,
    colback=blue!3,
    colframe=blue!50!black,
    boxrule=0.8pt,
    arc=2mm
}

\newtcolorbox{important}{
    enhanced,
    breakable,
    title={Importante},
    fonttitle=\bfseries,
    colback=yellow!8,
    colframe=orange!70!black,
    boxrule=0.8pt,
    arc=2mm
}

\newtcolorbox{example}[1]{
    enhanced,
    breakable,
    title={Ejemplo: #1},
    fonttitle=\bfseries,
    colback=green!4,
    colframe=green!40!black,
    boxrule=0.8pt,
    arc=2mm
}

\newtcolorbox{formula}[1]{
    enhanced,
    breakable,
    title={Fórmula / Regla: #1},
    fonttitle=\bfseries,
    colback=purple!4,
    colframe=purple!50!black,
    boxrule=0.8pt,
    arc=2mm
}

\newtcolorbox{exercise}[1]{
    enhanced,
    breakable,
    title={Ejercicio: #1},
    fonttitle=\bfseries,
    colback=gray!5,
    colframe=gray!60!black,
    boxrule=0.8pt,
    arc=2mm
}

\newtcolorbox{note}{
    enhanced,
    breakable,
    title={Nota},
    fonttitle=\bfseries,
    colback=white,
    colframe=gray!50,
    boxrule=0.6pt,
    arc=2mm
}

% ============================================================
% CONFIGURACIÓN FINAL
% ============================================================
\newcommand{\cornellhead}{%
\toprule
\textbf{Preguntas / palabras clave} & \textbf{Notas / respuestas} \\
\midrule
}

% ============================================================
% DOCUMENT
% ============================================================
\begin{document}

% ---------------- PORTADA ----------------
\begin{titlepage}
\thispagestyle{empty}

\begin{center}

\vspace*{1cm}

{\large \textsc{\universidad}}

\vspace{3cm}

{\Huge \textbf{\materia}}

\vspace{1cm}

{\LARGE \titulo}

\vspace{0.8cm}

\textit{Comprender los conceptos es el primer paso para convertir el estudio en conocimiento.}

\vspace{1.2cm}

\rule{0.70\textwidth}{0.5pt}

\vspace{0.5cm}

{\large Apuntes de estudio}

\vspace{0.5cm}

\rule{0.70\textwidth}{0.5pt}

\vfill

{\large \autor}

\vspace{0.3cm}

{\large \anio}

\end{center}

\vfill

\begin{flushleft}
\textit{Este es un modesto aporte a los estudiantes de ``Derecho Civil --- Personas Jurídicas'' no pretende sustituir el material oficial de la materia.}
\end{flushleft}

\end{titlepage}

% ---------------- ÍNDICE ----------------
\tableofcontents

% ============================================================
% CONTENIDO
% ============================================================

\chapter*{Presentación general}
\addcontentsline{toc}{chapter}{Presentación general}

Estos apuntes abordan de manera sistemática el régimen jurídico de las personas jurídicas en el Código Civil y Comercial de la Nación Argentina (Ley 26.994). Se analizan cuatro grandes ejes temáticos: (1) concepto, naturaleza jurídica, clasificación y atributos de las personas jurídicas; (2) regulación de las personas jurídicas privadas (nombre, domicilio, capacidad, patrimonio, estatutos y gobierno); (3) responsabilidad civil de las personas jurídicas y su relación con los administradores y dependientes; (4) fin de la existencia y tipos especiales: fundaciones, asociaciones civiles y simples asociaciones.

\begin{note}
Todo abogado, trabaje en derecho de familia, comercial, laboral, notarial o administrativo, se encuentra permanentemente con personas jurídicas. Saber identificar su tipo, leer su estatuto, comprender sus órganos de gobierno y conocer las reglas de responsabilidad es una habilidad transversal y cotidiana: desde redactar el estatuto de una ONG, hasta asesorar sobre la responsabilidad de los directivos de una sociedad anónima, gestionar la disolución de un consorcio o defender a una asociación civil en un proceso disciplinario.

Este conocimiento también resulta útil en la vida cotidiana, dado que las personas se relacionan constantemente con personas jurídicas: el consorcio del edificio, el club al que se pertenece, la empresa en la que se trabaja, el Estado que regula.
\end{note}

\textbf{Ejemplo integrador que guía la exposición:} un grupo de vecinos forma un club de fútbol. Para ello necesitan constituir una persona jurídica. ¿Qué tipo eligen? ¿Qué nombre pueden usar? ¿Quién representa al club? ¿Pueden sancionar a un socio? ¿Qué pasa si el club se disuelve? Cada bloque temático de estos apuntes responde a una de estas preguntas.

\section*{Mapa de contenidos}
\addcontentsline{toc}{section}{Mapa de contenidos}

La siguiente tabla resume los cuatro bloques temáticos en que se organiza la materia, con la pregunta guía correspondiente a cada punto.

\begin{longtable}{
    >{\raggedright\arraybackslash}p{0.10\textwidth}
    >{\raggedright\arraybackslash}p{0.28\textwidth}
    >{\raggedright\arraybackslash}X
}
\caption{Mapa de contenidos: Bloque I --- Teoría general de las personas jurídicas} \\
\toprule
\textbf{Punto} & \textbf{Tema} & \textbf{Pregunta guía} \\
\midrule
\endfirsthead
\toprule
\textbf{Punto} & \textbf{Tema} & \textbf{Pregunta guía} \\
\midrule
\endhead
1.1 & Concepto y definición (art. 141 CCyCN) & ¿Qué es una persona jurídica y quién le otorga ese carácter? \\
1.2 & Naturaleza jurídica --- teorías & ¿Existen realmente las personas jurídicas o son una creación del legislador? \\
1.3 & Clasificación: públicas y privadas (arts. 146, 147, 148) & ¿Cuáles son las diferencias entre una persona jurídica pública y una privada? \\
1.4 & Ley aplicable (art. 150) & ¿Qué norma prevalece cuando hay conflicto entre el estatuto y la ley? \\
1.5 & Personalidad diferenciada e inoponibilidad (art. 144) & ¿En qué casos se puede ignorar la personería y demandar directamente a los socios? \\
\bottomrule
\end{longtable}

\begin{longtable}{
    >{\raggedright\arraybackslash}p{0.10\textwidth}
    >{\raggedright\arraybackslash}p{0.28\textwidth}
    >{\raggedright\arraybackslash}X
}
\caption{Mapa de contenidos: Bloque II --- Atributos y funcionamiento} \\
\toprule
\textbf{Punto} & \textbf{Tema} & \textbf{Pregunta guía} \\
\midrule
\endfirsthead
\toprule
\textbf{Punto} & \textbf{Tema} & \textbf{Pregunta guía} \\
\midrule
\endhead
2.1 & Nombre (art. 151) & ¿Qué requisitos debe cumplir el nombre de una persona jurídica? \\
2.2 & Domicilio y sede social (arts. 152, 153) & ¿Cuál es la diferencia entre domicilio y sede social? \\
2.3 & Capacidad y principio de especialidad (arts. 141, 156) & ¿Puede una persona jurídica hacer cualquier cosa? \\
2.4 & Patrimonio (arts. 143, 154) & ¿Por qué el patrimonio de la persona jurídica es independiente al de sus socios? \\
2.5 & Estatutos y gobierno (arts. 158, 159, 160) & ¿Por qué es fundamental leer el estatuto antes de asesorar a una persona jurídica? \\
2.6 & Comienzo de la existencia (art. 142) & ¿Desde cuándo existe jurídicamente una persona jurídica? \\
\bottomrule
\end{longtable}

\begin{longtable}{
    >{\raggedright\arraybackslash}p{0.10\textwidth}
    >{\raggedright\arraybackslash}p{0.28\textwidth}
    >{\raggedright\arraybackslash}X
}
\caption{Mapa de contenidos: Bloque III --- Responsabilidad civil} \\
\toprule
\textbf{Punto} & \textbf{Tema} & \textbf{Pregunta guía} \\
\midrule
\endfirsthead
\toprule
\textbf{Punto} & \textbf{Tema} & \textbf{Pregunta guía} \\
\midrule
\endhead
3.1 & Evolución histórica --- art. 43 Código de Vélez & ¿Por qué el sistema anterior era considerado injusto? \\
3.2 & Responsabilidad por actos de directivos --- art. 1763 & ¿Cuándo responde la empresa por lo que hace su presidente? \\
3.3 & Responsabilidad por cosas y actividades riesgosas --- art. 1757 & ¿Responde la empresa aunque no haya culpa de nadie? \\
3.4 & Responsabilidad por dependientes --- art. 1753 & ¿Qué pasa si un empleado causa daño mientras trabaja? \\
3.5 & Responsabilidad de los administradores --- art. 160 & ¿Pueden los administradores ser demandados por la propia persona jurídica? \\
\bottomrule
\end{longtable}

\begin{longtable}{
    >{\raggedright\arraybackslash}p{0.10\textwidth}
    >{\raggedright\arraybackslash}p{0.28\textwidth}
    >{\raggedright\arraybackslash}X
}
\caption{Mapa de contenidos: Bloque IV --- Fin de la existencia y tipos especiales} \\
\toprule
\textbf{Punto} & \textbf{Tema} & \textbf{Pregunta guía} \\
\midrule
\endfirsthead
\toprule
\textbf{Punto} & \textbf{Tema} & \textbf{Pregunta guía} \\
\midrule
\endhead
4.1 & Causales de disolución --- art. 163 & ¿Cuándo y cómo puede terminar la existencia de una persona jurídica? \\
4.2 & Revocación estatal --- art. 164 & ¿Puede el Estado disolver una asociación civil arbitrariamente? \\
4.3 & Reconducción y liquidación & ¿Qué pasa con los bienes cuando se disuelve una persona jurídica? \\
4.4 & Fundaciones --- arts. 193--224 & ¿Qué diferencia a una fundación de una asociación civil? \\
4.5 & Asociaciones civiles --- arts. 168--186 & ¿Puede un club expulsar a un socio sin proceso previo? \\
4.6 & Simples asociaciones --- arts. 187--192 & ¿Cuál es la diferencia entre una simple asociación y una asociación civil? \\
\bottomrule
\end{longtable}

% ============================================================
\chapter{Introducción a las personas jurídicas}
% ============================================================

\section{Presentación del tema}

El estudio del régimen de las personas jurídicas en el derecho civil y comercial se sitúa dentro de la Parte General del Código, que se ocupa de los elementos de la relación jurídica: el sujeto, el objeto y la causa. Habiendo estudiado a fondo la persona humana, corresponde ahora tratar la persona jurídica, es decir, los entes colectivos que entablan también relaciones jurídicas en el campo del derecho privado.

\section{El problema de fondo: la asociatividad humana}

El problema central que resuelve esta regulación es cómo dar tratamiento jurídico a la asociatividad, dado que el ser humano es un ser sociable por naturaleza que tiende a unirse con otros para fines útiles. Esa regulación jurídica es fundamental porque permite establecer con qué requisitos esos entes colectivos, morales o ideales pueden entablar relaciones jurídicas: qué tipos de personas jurídicas existen, cuándo hay personalidad, cuáles son sus atributos, cuándo empieza y cuándo finaliza su existencia, y qué consecuencias tiene.

Estas personas jurídicas apuntan muchas veces a facilitar la organización de la vida cotidiana: desde el Estado como organización política que trasciende a los ciudadanos y tiene una personalidad distinta a la de estos, hasta las sociedades, las mutuales, las cooperativas y las iglesias.

\begin{note}
Existe un tipo de persona jurídica de enorme trascendencia económica, las sociedades, que tiene su materia específica (el Derecho Societario); el Estado también tiene la suya (el Derecho Administrativo). En esta materia se estudian las grandes disposiciones generales del Código Civil y Comercial aplicables a todas las personas jurídicas.
\end{note}

\section{Regulación en el Código de Vélez (Ley 340, 1871)}

\begin{note}
\textbf{Código de Vélez --- artículos relevantes.}

\textbf{Art. 30:} ``Todos los entes susceptibles de adquirir derechos y contraer obligaciones son personas''.

\textbf{Art. 31:} Las personas se clasifican en personas físicas (o de existencia visible) y personas jurídicas (o de existencia ideal).

\textbf{Art. 32:} Las personas jurídicas son definidas por exclusión: son todos los entes susceptibles de adquirir derechos y contraer obligaciones que no son personas de existencia visible.

\textbf{Arts. 33 a 50:} Régimen general de las personas jurídicas.
\end{note}

\section{El nuevo Código Civil y Comercial: metodología}

El nuevo Código regula la materia de forma más organizada y sistemática. El Libro Primero, la Parte General, comienza con un Título Primero sobre la persona humana, y el Título Segundo trata de las personas jurídicas.

Dentro de los seis libros del Código, en el Libro Primero:
\begin{itemize}
    \item Persona humana: artículos 19 al 140.
    \item Personas jurídicas: a partir del artículo 141.
\end{itemize}

El Título Segundo sobre personas jurídicas se divide en:
\begin{itemize}
    \item Capítulo 1: Parte general (personalidad, composición y clasificación).
    \item Capítulo 2: Asociaciones civiles y simples asociaciones.
    \item Capítulo 3: Fundaciones.
\end{itemize}

\begin{note}
El nuevo Código adopta la denominación única de ``persona jurídica'', diferenciándose del Código anterior, que hablaba de ``persona de existencia ideal'' y ``persona jurídica''. Antes de la Ley 17.711 (1968) se discutía si el género era la persona de existencia ideal y la persona jurídica una especie; con dicha reforma se tendió a equipararlos. Otras denominaciones usadas en doctrina son ``personas morales'' o ``personas colectivas''.
\end{note}

\section{Definición del artículo 141 del CCyCN}

\begin{definition}{Persona jurídica (art. 141 CCyCN)}
Son personas jurídicas todos los entes a los cuales el ordenamiento jurídico les confiere aptitud para adquirir derechos y contraer obligaciones para el cumplimiento de su objeto y los fines de su creación.
\end{definition}

Esta definición presenta dos características destacables:

\begin{enumerate}
    \item \textbf{Se define por la positiva}: se incorpora como elemento al ordenamiento jurídico, que es quien confiere la aptitud. Esto marca la diferencia con las personas humanas, cuya personalidad es un dato de la naturaleza reconocido —no concedido— por el legislador.
    \item \textbf{Se incorpora el elemento finalista}: la aptitud está orientada al cumplimiento del objeto y los fines de la creación. Este es el fundamento del principio de especialidad.
\end{enumerate}

\section{Principio de especialidad}

A diferencia de la persona humana, que puede dedicarse a cualquier actividad lícita, las personas jurídicas tienen un objeto y un fin propios. El \textbf{principio de especialidad} significa que la persona jurídica, dentro de su posibilidad de adquirir derechos y contraer obligaciones, está acotada a sus fines, a su especialidad, a su objeto, y a aquello para lo cual fue constituida.

\begin{important}
\textbf{Artículo 156 CCyCN --- Objeto.} El objeto de la persona jurídica privada debe ser preciso y determinado.
\end{important}

Esto resulta especialmente relevante para la redacción de los estatutos: los constituyentes deben ser claros, explícitos y descriptivos sobre cómo piensan lograr sus fines y cuál es el objeto de la persona jurídica.

\section{Teorías sobre la naturaleza jurídica}

\subsection{Teoría de la ficción (Savigny)}

Savigny recoge elaboraciones que provenían del derecho canónico y desarrolla la idea de la ficción. Para esta teoría, la persona humana es la única dotada de voluntad; como para Savigny el derecho es una expresión de la voluntad, no existirían propiamente personas jurídicas, sino un artificio creado por el legislador, que puede ser ampliado o restringido, pero que carece de voluntad propia.

\begin{important}
Consecuencia fundamental de la teoría de la ficción: la persona jurídica no respondería por los daños causados —responderían solo las personas humanas—, y como se trataría de una representación con mandato, este mandato nunca podría ser conferido para el ilícito. Esto dio lugar a situaciones injustas y, con el tiempo, se modificó esta posición. El artículo 43 del Código de Vélez, en su redacción original, sostenía esta posición: las personas jurídicas no respondían por los daños causados por quienes las dirigían o administraban.
\end{important}

\subsection{Teorías negatorias de la personalidad}

Niegan que existan propiamente personas jurídicas. Sostienen que lo que existe son patrimonios de afectación o bienes colectivos, pero no algo que pueda llamarse persona jurídica.

\subsection{Teorías de la realidad}

Buscan encontrar un sustrato real detrás de lo que se llama persona jurídica:

\begin{itemize}
    \item \textbf{Teoría organicista}: la persona jurídica sería como un organismo vivo, con sus mecanismos de adopción de la voluntad (los órganos que la gobiernan). Existe una realidad que el legislador viene a tutelar, proteger y regular.
    \item \textbf{Teoría de la institución} (Hauriou, seguida en Argentina por Llambías): las personas jurídicas son instituciones. Una institución requiere tres elementos: (a) una idea, objeto o fines en torno a los cuales se nuclean personas; (b) un sustrato personal, es decir, las personas que se nuclean en torno a esa idea; y (c) un poder u organización que permite llevar adelante esos fines.
\end{itemize}

\subsection{Teorías normativistas (Kelsen)}

Las personas jurídicas son un centro de imputación de normas: una realidad a la que el derecho le asigna consecuencias jurídicas a través de normas dictadas por la autoridad competente.

\begin{note}
En el Código Civil y Comercial vigente se percibe una fuerte presencia de las teorías normativistas —el ordenamiento jurídico es quien confiere la aptitud—, pero también existe una dimensión institucional, en tanto se reconocen realidades organizativas como las simples asociaciones.
\end{note}

\section{Clasificación: personas jurídicas públicas y privadas}

El Código Civil y Comercial simplifica la clasificación entre públicas y privadas.

\subsection{Personas jurídicas públicas (arts. 146 y 147)}

\begin{definition}{Personas jurídicas públicas (art. 146 CCyCN)}
Son personas jurídicas públicas:
\begin{enumerate}[label=\alph*)]
    \item El Estado nacional, las provincias, la Ciudad Autónoma de Buenos Aires, los municipios, las entidades autárquicas y las demás organizaciones a las que el ordenamiento jurídico les atribuya ese carácter.
    \item Los estados extranjeros, las organizaciones a las que el derecho internacional público reconozca personalidad jurídica (como la ONU) y toda otra persona jurídica constituida en el extranjero cuyo carácter público resulte de su derecho aplicable.
    \item La Iglesia Católica.
\end{enumerate}
\end{definition}

La inclusión de la Iglesia Católica con personería jurídica pública se vincula con el artículo 2 de la Constitución Nacional, que establece que el gobierno federal sostiene el culto católico apostólico romano, en virtud de la preeminencia histórica de la Iglesia antes de la constitución del Estado argentino.

En cuanto a su reconocimiento, comienzo, capacidad, funcionamiento y organización, las personas jurídicas públicas se rigen por sus propias leyes y ordenamientos. Así, por ejemplo, si se litiga contra el Banco Central, habrá que revisar su Carta Orgánica y las leyes que lo regulan.

\subsection{Personas jurídicas privadas (art. 148)}

El artículo 148 enumera los tipos de personas jurídicas privadas.

\begin{table}[H]
\centering
\caption{Tipos de personas jurídicas privadas (art. 148 CCyCN)}
\begin{tabularx}{\textwidth}{
    >{\raggedright\arraybackslash}p{0.22\textwidth}
    >{\raggedright\arraybackslash}X
    >{\raggedright\arraybackslash}p{0.22\textwidth}
}
\toprule
\textbf{Tipo} & \textbf{Característica principal} & \textbf{Ley aplicable} \\
\midrule
Sociedades & Finalidad directa de lucro. Los socios participan en ganancias y pérdidas. & Ley 19.550 \\
Asociaciones civiles & No tienen fin de lucro. Unión de personas para un fin común (cultural, deportivo, social, etc.). & Arts. 168--186 CCyCN \\
Simples asociaciones & Similar a la asociación civil, pero sin autorización estatal para funcionar. & Arts. 187--192 CCyCN \\
Fundaciones & Patrimonio afectado por uno o varios fundadores a un fin de bien común. Sin miembros. & Arts. 193--224 CCyCN \\
Iglesias y comunidades religiosas (no católicas) & No tienen fin de lucro. Libertad de culto y conciencia. & Ley 21.745 \\
Mutuales & Ayuda mutua entre sus miembros. Socorros recíprocos. & Ley 20.321 \\
Cooperativas & Asociación de personas para fines comunes de producción, trabajo o consumo. & Ley 20.337 \\
Consorcios de propiedad horizontal & Conjunto de propietarios de unidades funcionales de un inmueble sometido al régimen de propiedad horizontal. & Art. 2044 CCyCN \\
\bottomrule
\end{tabularx}
\end{table}

\begin{note}
El artículo 149 establece que, si el Estado participa en una persona jurídica privada, ello no modifica el carácter privado de esta, aunque puede haber derechos y obligaciones diferenciadas en razón del interés público comprometido.
\end{note}

\section{Ley aplicable (art. 150)}

\begin{formula}{Ley aplicable a las personas jurídicas privadas (art. 150 CCyCN)}
Las personas jurídicas privadas se rigen:
\begin{enumerate}
    \item[1.°)] Por las normas imperativas de la ley especial o, en su defecto, de este Código;
    \item[2.°)] Por las normas del acto constitutivo con sus modificaciones y de los reglamentos, prevaleciendo las primeras en caso de divergencia;
    \item[3.°)] Por las normas supletorias de leyes especiales, o en su defecto, por las de este Título.
\end{enumerate}
\end{formula}

El estatuto tiene fuerza de norma jurídica. Por ello, siempre que se aborde una cuestión en la que interviene una persona jurídica, lo primero que corresponde hacer es leer su estatuto en su última versión, con todas las modificaciones, y también sus reglamentos.

\section{Personalidad diferenciada e inoponibilidad (art. 144)}

Un presupuesto fundamental es que la persona jurídica es distinta de sus miembros. Esto implica distintos patrimonios, distintas responsabilidades, y la posibilidad incluso de que el propio miembro demande a la persona jurídica o viceversa.

\begin{important}
\textbf{Regla general}: los miembros no responden por las obligaciones de la persona jurídica. Esta es una de las diferencias más importantes que debe considerar quien constituye una persona jurídica: responde con el capital que aporta, pero no con todo su patrimonio personal.
\end{important}

\begin{definition}{Inoponibilidad de la persona jurídica (art. 144 CCyCN)}
La actuación que esté destinada a la consecución de fines ajenos a la persona jurídica, que constituya un recurso para violar la ley, el orden público o la buena fe, o para frustrar derechos de cualquier persona, se imputa a quienes, a título de socios, asociados, miembros o controlantes directos o indirectos, la hicieron posible, quienes responderán solidaria e ilimitadamente por los perjuicios causados.
\end{definition}

Esta norma reconoce como antecedente el artículo 54 de la Ley de Sociedades (Ley 19.550), pero ahora se aplica a todo el sistema de personas jurídicas por estar ubicada en la Parte General.

La inoponibilidad significa que los efectos de la personería no pueden hacerse valer frente a quien ha sido dañado, quien puede accionar directamente contra los socios, asociados, miembros o controlantes que hicieron posible esa actuación dañina. Responden solidaria e ilimitadamente: todos responden por el todo y con todos sus patrimonios.

% ============================================================
\chapter{Régimen de las personas jurídicas privadas}
% ============================================================

El Código ofrece reglas generales sobre nombre, domicilio, capacidad, patrimonio, estatutos y comienzo de la existencia, que luego cada tipo de persona jurídica complementa con sus normas específicas.

\section{Nombre (art. 151)}

\begin{definition}{Nombre de la persona jurídica (art. 151 CCyCN)}
La persona jurídica debe tener un nombre que la identifique como tal, con el aditamento indicativo de la forma jurídica adoptada. La persona jurídica en liquidación debe aclarar esta circunstancia en la utilización de su nombre. El nombre debe satisfacer recaudos de veracidad, novedad y aptitud distintiva, tanto respecto de otros nombres como de marcas, nombres de fantasía u otras formas de referencia a bienes o servicios, se relacionen o no con el objeto de la persona jurídica.

No puede contener términos o expresiones contrarios a la ley, el orden público o las buenas costumbres, ni inducir a error sobre la clase u objeto de la persona jurídica.

Si se incluye el nombre de personas humanas, deben contar con la conformidad de estas, que se presume si son socios. Sus herederos pueden oponerse a la continuación del uso si acreditan perjuicios materiales o morales.
\end{definition}

Los requisitos pueden precisarse del siguiente modo:

\begin{itemize}
    \item \textbf{Indicación de la forma jurídica}: debe agregarse el tipo de persona jurídica (por ejemplo, ``El estudiante feliz S.A.''). Si está en liquidación, se agrega ``en liquidación''.
    \item \textbf{Veracidad}: concordancia entre el nombre y el objeto, para que no haya lugar a confusión. Esto se juega especialmente en ciertos ámbitos, como el de las universidades.
    \item \textbf{Novedad y aptitud distintiva}: el nombre debe distinguirse de otros nombres, marcas, nombres de fantasía u otras formas de referencia a bienes y servicios.
\end{itemize}

\begin{example}{Teatro Colón S.A.}
No se puede fundar una sociedad anónima con el nombre ``Teatro Colón S.A.'' porque se confundiría con el Teatro Colón preexistente. El nombre no tendría novedad ni aptitud distintiva.
\end{example}

\begin{example}{Fundación Lionel Messi / Fundación Borges}
Si se quiere usar el nombre de una persona humana, se necesita su conformidad. Si esa persona ya falleció, los herederos pueden oponerse si acreditan perjuicios materiales o morales. Así, si la Fundación Borges —supuestamente dedicada a las altas letras— organizara un concurso de letras de reggaetón, los herederos de Borges podrían oponerse alegando que ello afecta el buen nombre del escritor y resulta incompatible con los fines fundacionales.
\end{example}

\section{Domicilio y sede social (arts. 152 y 153)}

El domicilio es el asiento jurídico de la persona jurídica y permite ubicarla, determinar el juez competente y conocer su inscripción registral.

\begin{table}[H]
\centering
\caption{Domicilio y sede social}
\begin{tabularx}{\textwidth}{>{\raggedright\arraybackslash}X >{\raggedright\arraybackslash}X}
\toprule
\textbf{Domicilio} & \textbf{Sede social} \\
\midrule
Se fija en el estatuto. Puede expresarse simplemente con una ciudad (por ejemplo, ``Buenos Aires''). & Es la dirección física concreta donde tienen lugar la dirección y administración de la persona jurídica (por ejemplo, ``Av. Figueroa Alcorta 2263, CABA''). \\
Determina el juez competente y la jurisdicción registral. & Generalmente se inscribe ante el organismo de contralor (IGJ). Si está en el estatuto, cambiarla requiere modificar el estatuto. \\
\bottomrule
\end{tabularx}
\end{table}

\begin{important}
\textbf{Artículo 153 CCyCN --- Efectos de la sede social.} Se tienen por válidas y vinculantes para la persona jurídica todas las notificaciones efectuadas en la sede inscripta.
\end{important}

También existe el domicilio especial: cuando la persona jurídica posee establecimientos o sucursales, tiene su domicilio especial en el lugar de dichos establecimientos, pero solo para las obligaciones allí contraídas.

\begin{example}{Sucursal del Banco Santander en Mendoza}
Si un contratista de pintura firmó un contrato con la sucursal mendocina del Banco Santander y este no le paga, puede litigar en Mendoza. No necesita ir a Buenos Aires a litigar contra la casa matriz, porque la obligación se contrajo localmente.
\end{example}

\section{Capacidad y principio de especialidad (arts. 141, 156)}

La persona jurídica es capaz de derecho; la capacidad de ejercicio se realiza a través de sus órganos representativos. Pero esa capacidad tiene límites:

\begin{enumerate}
    \item \textbf{Principio de especialidad}: la capacidad está acotada al objeto y a los fines de la persona jurídica.
    \item \textbf{Limitaciones legales}: la ley prohíbe a las personas jurídicas adquirir ciertos derechos o los limita. Por ejemplo, el usufructo en favor de personas jurídicas tiene un plazo determinado, mientras que para personas humanas puede extenderse hasta su muerte.
    \item \textbf{Limitaciones por la naturaleza de las cosas}: las personas jurídicas no pueden adoptar niños, porque el derecho de familia se refiere a las personas humanas.
\end{enumerate}

\section{Patrimonio (arts. 143 y 154)}

La persona jurídica debe tener un patrimonio distinto del de sus miembros. Esto es una consecuencia directa de la personalidad diferenciada.

La persona jurídica tiene bienes a su nombre; si se trata de información registrable, se pueden inscribir preventivamente. Por ejemplo, si alguien dona un inmueble a una fundación y esta tarda en inscribirse, puede realizarse una anotación preventiva en el registro para proteger el bien mientras se completan los trámites.

\begin{note}
En cuanto a su duración, las personas jurídicas son en principio ilimitadas en el tiempo, excepto que el acto constitutivo disponga lo contrario. La prórroga está regulada en el artículo 165 y requiere decisión de los miembros o aprobación de la autoridad de contralor antes del vencimiento del plazo.
\end{note}

\section{Comienzo de la existencia (art. 142)}

\begin{definition}{Comienzo de la existencia (art. 142 CCyCN)}
La existencia de la persona jurídica privada comienza desde su constitución. No necesita autorización legal para funcionar, excepto disposición legal en contrario. En los casos en que se requiere autorización estatal, la persona jurídica no puede funcionar antes de obtenerla.
\end{definition}

En el Código anterior existía una discusión sobre si el comienzo de la existencia era el momento de la constitución o el de la autorización. El nuevo Código es claro: el comienzo es la constitución. Sin embargo, si la ley requiere autorización estatal para funcionar (como en el caso de las fundaciones), la persona jurídica no puede comenzar a operar hasta obtenerla.

\begin{example}{Fundación banco de alimentos}
Si hoy se constituye una fundación con objeto de banco de alimentos y se aporta un capital de \$100.000, la fundación existe desde ese momento. Pero no puede empezar a hacer compras, emitir facturas ni hacer donaciones hasta que se obtenga la autorización estatal. La constitución y la autorización son dos momentos distintos.
\end{example}

\section{Estatutos y gobierno (arts. 150, 158, 159, 160)}

El estatuto es una norma jurídica secundaria que emana de las partes que constituyen la persona jurídica y que la regula internamente. No debe confundirse con el acto constitutivo:

\begin{itemize}
    \item El \textbf{acto constitutivo} es el acto jurídico por el cual las partes deciden dar inicio a la persona jurídica. Generalmente, en ese acto se redacta el estatuto.
    \item El \textbf{estatuto} cobra independencia porque puede sufrir sucesivas modificaciones a lo largo del tiempo.
\end{itemize}

\begin{note}
Hay clubes que tienen más de cien años y cuyo estatuto ha cambiado muchísimas veces. Cuando se solicita el estatuto ante la Inspección de Justicia, se pide el actualizado con todos sus cambios.
\end{note}

\subsection{Contenido del estatuto}

El estatuto de una persona jurídica privada debe contener:

\begin{itemize}
    \item Nombre y domicilio.
    \item Objeto preciso y determinado (art. 156).
    \item Derechos y deberes de los miembros.
    \item Forma de gobierno: composición de los órganos, cómo se toman las decisiones, quórum y mayorías necesarias (incluidas las agravadas para decisiones trascendentes, como reforma del estatuto, expulsión de miembros o disolución).
    \item Atribuciones de cada órgano y sus miembros.
    \item Administración de bienes.
    \item Destino de los bienes en caso de disolución.
\end{itemize}

\subsection{Deber de lealtad y diligencia de los administradores (art. 159)}

\begin{important}
\textbf{Artículo 159 CCyCN --- Deber de lealtad y diligencia.} Los administradores de la persona jurídica deben obrar con lealtad y diligencia. No pueden perseguir ni favorecer intereses contrarios a los de la persona jurídica. Si en determinada operación tuvieran un interés contrario al de ella, deben hacerlo saber a los demás miembros del órgano de administración o, en su caso, al órgano de gobierno, y abstenerse de cualquier intervención relacionada con dicha operación.
\end{important}

Ser administrador de una persona jurídica implica una gran responsabilidad, porque se confían bienes económicos que no son propios.

% ============================================================
\chapter{Responsabilidad civil de las personas jurídicas}
% ============================================================

\section{Evolución histórica: del artículo 43 de Vélez al nuevo Código}

El Código de Vélez Sarsfield, adherente a la teoría de la ficción, establecía en el artículo 43 que las personas jurídicas no respondían por los daños que causaban quienes las dirigían o administraban: no se podía iniciar una acción de daños y perjuicios contra una persona jurídica.

La lógica era coherente con la teoría de la ficción: como la personería era solo ficción y el mandato nunca se confería para el ilícito, las actuaciones ilícitas se imputaban a las personas humanas, no a la persona jurídica. A lo largo del tiempo, esta solución fue generando la sensación de que constituía una injusticia.

% TODO: contenido incompleto o ambiguo en las notas originales — la referencia al profesor "Rivarola" y a la traducción del término utilizado por Freitas se transcribe tal como figura en el apunte, sin poder verificarse contra la fuente original.
Existieron distintas posturas doctrinarias para resolver el problema. Un profesor de apellido Rivarola acudió al diccionario que empleaba Vélez y encontró que, en el texto original de Freitas (en portugués), el término utilizado podía entenderse como ``cuando'' y no como ``aunque'', lo que cambiaría el sentido de la norma; sin embargo, esta interpretación no fue pacífica. El jurista Borda, en un fallo, llegó a declarar que el artículo 43 era de una injusticia intolerable.

\begin{important}
Todo cambió en 1968 con la Ley 17.711, que reformó el artículo 43 y estableció que la persona jurídica responde por los daños que causen quienes la dirigen o administran en ejercicio o con ocasión de sus funciones. El nuevo Código recoge esta solución en el artículo 1763.
\end{important}

\section{Tres supuestos de responsabilidad}

\subsection{Responsabilidad por actos de directivos y administradores (art. 1763)}

\begin{formula}{Responsabilidad de la persona jurídica (art. 1763 CCyCN)}
La persona jurídica responde por los daños que causen quienes las dirigen o administran en ejercicio o con ocasión de sus funciones.
\end{formula}

La expresión ``en ejercicio o con ocasión de sus funciones'' resulta determinante: si un presidente de la empresa está de vacaciones y comete un daño, la empresa no responde, porque el acto no ocurre en ejercicio ni con ocasión de sus funciones.

\subsection{Responsabilidad por cosas y actividades riesgosas (art. 1757)}

\begin{formula}{Hecho de las cosas y actividades riesgosas (art. 1757 CCyCN)}
Toda persona responde por el daño causado por el riesgo o vicio de las cosas, o de las actividades que sean riesgosas o peligrosas por su naturaleza, por los medios empleados o por las circunstancias de su realización.
\end{formula}

Esta norma se aplica a todas las personas, incluidas las jurídicas. Se trata de una responsabilidad objetiva: no depende de factores de atribución subjetivos (dolo o culpa); la responsabilidad se atribuye siempre.

\begin{example}{Sociedad anónima y accidente de tránsito}
Una sociedad anónima tiene un automóvil. Ese auto causa un daño. La persona jurídica responde, aunque no estuviera manejando: se le imputa la responsabilidad porque es dueña de la cosa riesgosa.
\end{example}

\subsection{Responsabilidad por los dependientes (art. 1753)}

\begin{formula}{Responsabilidad del principal por el hecho del dependiente (art. 1753 CCyCN)}
El principal responde objetivamente por los daños que causen los que están bajo su dependencia, o las personas de las cuales se sirve para el cumplimiento de sus obligaciones, cuando el hecho dañoso acaece en ejercicio o con ocasión de las funciones encomendadas.
\end{formula}

También esta norma engloba a las personas jurídicas. El término ``dependientes'' es genérico: incluye a cualquier persona de la que la persona jurídica se sirve para cumplir sus obligaciones, más allá de la relación de empleo formal.

\begin{example}{Personal de seguridad que golpea a alguien}
En una empresa, el personal de seguridad golpea a una persona de manera arbitraria y violenta. El daño causado se imputa tanto al empleado como a la persona jurídica. La empresa no puede escudarse en que ``fue el empleado''.
\end{example}

\section{Responsabilidad de los administradores (art. 160)}

\begin{formula}{Responsabilidad de los administradores (art. 160 CCyCN)}
Los administradores responden en forma ilimitada y solidaria frente a la persona jurídica, sus integrantes y terceros, por los daños causados por su culpa en el ejercicio o con ocasión de sus funciones, por acción u omisión.
\end{formula}

Esta disposición es la contracara del artículo 1763: si el presidente causó un daño, frente al tercero dañado responden tanto el presidente como la persona jurídica. Pero la persona jurídica, a su vez, puede iniciar acción contra el presidente por los daños que causó. Los administradores responden de forma ilimitada (con todos sus bienes) y solidaria (todos responden por el todo).

% ============================================================
\chapter{Fin de la existencia y fundaciones}
% ============================================================

\section{Causales de disolución (art. 163)}

\begin{formula}{Causales de disolución (art. 163 CCyCN)}
La persona jurídica se disuelve por:
\begin{enumerate}[label=\alph*)]
    \item La decisión de sus miembros adoptada por unanimidad o por la mayoría establecida por el estatuto o disposición especial;
    \item El cumplimiento de la condición resolutoria a la que el acto constitutivo subordinó su existencia;
    \item La consecución del objeto para el cual la persona jurídica se formó, o la imposibilidad sobreviniente de cumplirlo;
    \item El vencimiento del plazo;
    \item La declaración de quiebra;
    \item La fusión respecto de las personas jurídicas que se fusionan o la escisión respecto de la escindida;
    \item La reducción a uno del número de miembros, si la ley especial exige pluralidad de ellos y esta no es restablecida dentro de los tres meses;
    \item La denegatoria o revocación firmes de la autorización estatal para funcionar, cuando esta sea requerida;
    \item El agotamiento de los bienes destinados a sostenerla;
    \item Cualquier otra causa prevista en el estatuto o en disposiciones de este Título o de ley especial.
\end{enumerate}
\end{formula}

Entre las causales más relevantes cabe señalar:

\begin{itemize}
    \item \textbf{Inciso a) --- Decisión de los miembros}: por unanimidad si el estatuto no dice nada; si el estatuto establece una mayoría, se aplica esa mayoría.
    \item \textbf{Inciso b) --- Condición resolutoria}: cuando la existencia estaba subordinada a un hecho futuro e incierto que se cumple.
    \item \textbf{Inciso c) --- Consecución del objeto}: cuando la persona jurídica se creó para un fin específico y este se cumplió o se tornó imposible.
    \item \textbf{Inciso e) --- Quiebra}: cuando el pasivo supera al activo y no hay propuesta de acuerdo preventivo viable. La quiebra está regulada por ley especial y tiene por objeto liquidar el activo para pagar a los acreedores; excepcionalmente, puede haber avenimiento o conversión.
    \item \textbf{Inciso h) --- Revocación estatal}: el Estado puede revocar la autorización para funcionar, con las restricciones del artículo 164.
\end{itemize}

\begin{example}{Fundación condicionada a un programa estatal}
Una fundación se crea para ayudar a personas en situación de calle ``siempre y cuando no exista un programa especial del Gobierno de la Ciudad de Buenos Aires para esa problemática''. Si el gobierno implementa tal programa, se cumple la condición y la fundación se disuelve.
\end{example}

\begin{example}{Fundación para los cien años de un club}
Se constituye una fundación especial para organizar la celebración del centenario de un club. Pasada la celebración, el objeto se cumplió y la personería se disuelve.
\end{example}

\section{Revocación estatal --- garantías (art. 164)}

\begin{important}
\textbf{Artículo 164 CCyCN --- Revocación de la autorización estatal.} La revocación de la autorización estatal debe fundarse en la comisión de actos graves que violen la ley, el estatuto o el reglamento. La resolución que la disponga debe ser fundada y contener la indicación de los actos que la justifican. Se debe dar a los interesados la posibilidad de ser escuchados antes de que la resolución se dicte. Es apelable judicialmente. La apelación puede tener efectos suspensivos.
\end{important}

Las garantías que protegen a la persona jurídica frente a una revocación arbitraria pueden sintetizarse así:

\begin{enumerate}
    \item Solo procede por actos graves que violen la ley, el estatuto o el reglamento.
    \item La resolución debe ser fundada (no puede ser arbitraria).
    \item Se debe garantizar el derecho de defensa (descargo previo).
    \item Es apelable judicialmente, y la apelación puede tener efecto suspensivo (la persona jurídica sigue funcionando mientras se tramita).
\end{enumerate}

El fundamento constitucional de estas garantías es la libertad de asociación consagrada en el artículo 14 de la Constitución Nacional.

\section{Reconducción y liquidación}

\begin{itemize}
    \item \textbf{Reconducción}: si se decidió la disolución pero no terminó la liquidación, se puede revertir esa decisión siempre que la causa pueda ser removida. Por ejemplo, si se cumplió el plazo pero los miembros deciden ampliarlo y transformar la personería de tiempo limitado a indefinido.
    \item \textbf{Liquidación}: una vez dada la causal de disolución, se pagan las deudas y obligaciones pendientes con el activo, se abonan los gastos, y los bienes remanentes se destinan según lo establezca la ley o el estatuto. En asociaciones civiles y fundaciones generalmente se exige que los bienes vayan a otra institución de bien común, sin fin de lucro.
\end{itemize}

\section{Las fundaciones (arts. 193--224)}

\begin{definition}{Fundación (art. 193 CCyCN)}
Las fundaciones son personas jurídicas que se constituyen con una finalidad de bien común, sin propósito de lucro, mediante el aporte patrimonial de una o más personas, denominadas fundadores, destinando dicho aporte al cumplimiento de sus fines.
\end{definition}

Características distintivas de las fundaciones:

\begin{itemize}
    \item No tienen miembros: el fundador aporta un patrimonio, que queda afectado al cumplimiento de un fin de bien común.
    \item No tienen fin de lucro.
    \item Requieren autorización estatal para funcionar (art. 142).
    \item Pueden constituirse por acto entre vivos o por testamento.
\end{itemize}

\begin{example}{Fundación banco de alimentos (caso detallado)}
En un contexto de crisis, un fundador con capital decide constituir una fundación con \$100.000 para organizar un banco de alimentos. La fundación alquila un depósito, contrata empleados, establece convenios con instituciones que proveen alimentos y organiza campañas de donación. Todo eso es la fundación funcionando. No hay ``socios'': hay un fundador que dotó el patrimonio y un Consejo de Administración que lo gestiona.

También puede constituirse por testamento: un deportista que practique hockey puede establecer en su testamento que su parte disponible de la herencia vaya a una fundación para que personas sin recursos puedan practicar ese deporte.
\end{example}

\subsection{Patrimonio inicial suficiente}

El Código exige que la fundación cuente con un patrimonio inicial que posibilite razonablemente el cumplimiento de sus fines (recaudo no exigido para las asociaciones civiles). La Inspección de Personas Jurídicas no otorgará la autorización si los fondos son insuficientes.

El plan de acción a tres años (plan trienal) que se presenta ante el organismo de contralor puede incluir tanto los aportes iniciales como los ingresos futuros proyectados (eventos benéficos, campañas de donación, etc.).

\subsection{Órgano de gobierno: Consejo de Administración}

El único órgano de gobierno es el Consejo de Administración, dotado de todas las atribuciones para el gobierno y la administración. Debe tener al menos tres personas humanas como integrantes.

Puede delegar funciones en un Comité Ejecutivo (art. 205), que ejerce esas funciones entre las reuniones del Consejo. Los miembros del Comité Ejecutivo pueden ser remunerados, mientras que los del Consejo de Administración no.

\subsection{Control externo}

Las fundaciones no tienen órgano de fiscalización interna: el control es exclusivamente externo, a cargo de la autoridad de contralor (IGJ u organismo provincial equivalente). Entre las atribuciones del organismo de contralor (art. 222) se encuentran:

\begin{itemize}
    \item Aprobar estatutos y sus reformas.
    \item Fiscalizar el funcionamiento.
    \item Designar administradores interinos, suspender deliberaciones, remover administradores, convocar asambleas.
    \item Intervenir en la disolución y liquidación.
\end{itemize}

\subsection{Inmutabilidad del objeto (art. 216)}

No es posible reformar el objeto de la fundación, salvo que haya resultado de cumplimiento imposible. Esto se debe a que la fundación fue autorizada para ese objeto específico, y cambiarlo implicaría traicionar la voluntad inicial del fundador.

% ============================================================
\chapter{Asociaciones civiles y simples asociaciones}
% ============================================================

Las asociaciones civiles y las simples asociaciones son personas jurídicas privadas de enorme importancia social: clubes deportivos, ONG, asociaciones barriales, organizaciones de víctimas, entidades religiosas. La Constitución Nacional consagra la libertad de asociación como un derecho fundamental.

\section{Asociaciones civiles (arts. 168--186)}

\subsection{Concepto y objeto}

\begin{definition}{Objeto de las asociaciones civiles (art. 168 CCyCN)}
La asociación civil debe tener un objeto que no sea contrario al interés general o al bien común. El objeto no puede ser contrario a las disposiciones de este Código ni a la ley especial, ni referirse a cuestiones de interés particular de los asociados.
\end{definition}

La definición es por la negativa: cualquier finalidad que no afecte al interés general puede ser objeto de una asociación civil. La asociación civil es una corporación: unión de personas en torno a un fin común (cultural, artístico, deportivo, religioso, social).

\begin{important}
La asociación civil no tiene fin de lucro y no persigue el lucro como fin principal. Puede eventualmente generar ganancias si organiza un evento a beneficio, pero ese dinero se reinvierte en los fines de la asociación; no se distribuye entre los miembros.
\end{important}

\subsection{Diferencia fundamental con las sociedades}

En la sociedad, los socios aportan capital, participan de ganancias y pérdidas, y obtienen utilidades. En la asociación civil, los miembros aportan su cuota —fijada por el estatuto y las asambleas— que contribuye a los fines colectivos, sin distribuirse ganancias.

\begin{note}
La Resolución de la IGJ 7/2015 entiende por ``bien común'' a las finalidades que contribuyen al bien de la comunidad en general o a la mejora de las condiciones de vida, en contraposición al bien individual.
\end{note}

\subsection{Constitución y contenido del estatuto (art. 170)}

El instrumento constitutivo debe hacerse en escritura pública. El artículo 170 establece que debe contener:

\begin{itemize}
    \item Identificación de los constituyentes.
    \item Nombre, objeto y domicilio.
    \item Plazo de duración y causales de disolución.
    \item Composición del patrimonio inicial.
    \item Clases o categorías de socios (activos, cadetes, vitalicios, etc.).
    \item Régimen de admisión, renuncia y sanciones disciplinarias.
    \item Órganos de gobierno: composición, atribuciones, funcionamiento y duración de cargos.
    \item Cierre del ejercicio económico anual.
    \item Destino de los bienes en caso de liquidación: deben ir a otra entidad de bien común, pública o privada, sin fin de lucro, con domicilio en la República Argentina.
\end{itemize}

\section{Órganos de gobierno}

\subsection{Comisión Directiva (art. 171)}

Es el órgano de dirección y administración cotidiana de la asociación civil. Los cargos mínimos son Presidente, Secretario y Tesorero (los demás son Vocales). Cada estatuto puede agregar cargos adicionales: vicepresidentes primero y segundo, prosecretario, protesorero, etc.

Debe reunirse al menos una vez al mes. Sus atribuciones las fija el estatuto, y hay que distinguir entre las que tiene la Comisión como cuerpo colegiado y las que tiene el Presidente individualmente.

\begin{important}
No hay que confundir el hecho de que alguien sea presidente con que pueda resolver todo por sí mismo. El presidente es quien preside un órgano que es la Comisión Directiva.
\end{important}

\begin{example}{Venta de un inmueble por una asociación civil}
Si una asociación civil quiere vender un inmueble, hay que leer el estatuto para saber quién puede tomar esa decisión. A veces se exige que la asamblea autorice previamente la venta; otras veces basta con la Comisión Directiva. El abogado debe siempre revisar el estatuto antes de asesorar.
\end{example}

\subsection{Asamblea}

Es el órgano supremo: la reunión de todos los asociados. En asociaciones muy numerosas puede instrumentarse a través de representantes elegidos por los socios.

\begin{itemize}
    \item \textbf{Asamblea ordinaria}: generalmente una vez al año; aprueba la memoria y el balance presentado por la Comisión Directiva con el dictamen del órgano de fiscalización.
    \item \textbf{Asamblea extraordinaria}: se convoca cuando es necesario modificar el estatuto, elegir autoridades, aprobar reglamentos o tomar decisiones que requieran mayorías especiales.
\end{itemize}

\subsection{Órgano de fiscalización}

Puede recaer en personas no asociadas. Si la asociación cuenta con más de cien socios, debe tener una Comisión Revisora de Cuentas. Si la asociación exige una profesión determinada para ser socio, los integrantes del órgano de fiscalización no necesariamente deben contar con ese título, pero debe contratarse a profesionales independientes para el asesoramiento correspondiente.

Además del control interno, existe también la fiscalización estatal del artículo 174.

\section{Poder disciplinario}

La asociación civil puede sancionar a sus miembros. En esto se evidencia la distinción entre la persona jurídica y sus miembros.

\subsection{Principios que rigen el poder disciplinario}

\begin{enumerate}
    \item Debe estar regulado en el estatuto, con claridad sobre qué órgano interviene (asamblea, Comisión Directiva o Comisión de Disciplina).
    \item Se debe garantizar el derecho de defensa del afectado: debe darse vista al sumariado y permitirle hacer sus descargos.
    \item Gradualidad y proporcionalidad de la sanción: primero se amonesta, luego se suspende por períodos crecientes, y finalmente se expulsa. No se puede expulsar sin proceso.
    \item La decisión es apelable ante los tribunales judiciales si no se respetó el derecho de defensa, si fue desproporcionada o si no hubo debido proceso.
\end{enumerate}

\begin{example}{Destrucción de un vestuario}
Un socio muy enojado, luego de un partido, destruye el vestuario del club un domingo. El lunes, la Comisión Directiva puede suspenderlo transitoriamente como medida cautelar. Pero no puede expulsarlo sin proceso: debe abrirse un sumario, recopilarse pruebas y testimonios, notificarse al imputado, permitirle hacer su descargo y luego resolver. La sanción debe ser proporcional a la falta.
\end{example}

\subsection{Límites de la revisión judicial}

El juez puede revisar la forma del proceso disciplinario (si se respetó el derecho de defensa, si la sanción fue proporcional). Pero no puede revisar el fondo de la decisión —salvo grave violación de derechos humanos—, porque ello implicaría alterar el objeto y la autonomía de la persona jurídica.

\begin{example}{Club vegano y el vendedor de panchos}
Un club de personas veganas expulsa a un socio que insiste en vender panchos en las instalaciones. Si el juez le diera la razón al expulsado en cuanto al fondo, estaría alterando el objeto de la asociación y el derecho de todos sus miembros a mantener la identidad que los une. Por eso, el control judicial se centra en la forma, no en el fondo de la decisión disciplinaria.
\end{example}

\section{Simples asociaciones (arts. 187--192)}

Las simples asociaciones son personas jurídicas privadas que, a diferencia de las asociaciones civiles, no necesitan autorización estatal para funcionar. Existen muchísimas en la práctica.

\subsection{Diferencias con las asociaciones civiles}

\begin{table}[H]
\centering
\caption{Asociación civil frente a simple asociación}
\begin{tabularx}{\textwidth}{
    >{\raggedright\arraybackslash}p{0.24\textwidth}
    >{\raggedright\arraybackslash}X
    >{\raggedright\arraybackslash}X
}
\toprule
\textbf{Aspecto} & \textbf{Asociación civil} & \textbf{Simple asociación} \\
\midrule
Instrumento constitutivo & Escritura pública & Instrumento público o privado con firma certificada \\
Autorización estatal & Requerida & No requerida \\
Comienzo de existencia & Desde la constitución (previa autorización para funcionar) & Desde la constitución \\
Órgano de fiscalización & Obligatorio (Comisión Revisora de Cuentas si hay más de 100 socios) & Prescindible si hay menos de 20 socios; se certifica con contador externo \\
Responsabilidad en insolvencia & Separación plena entre persona jurídica y miembros & Administradores y miembros que participaron responden solidariamente con sus bienes propios \\
Régimen supletorio & No aplica & Se aplica el régimen de asociaciones civiles \\
\bottomrule
\end{tabularx}
\end{table}

\subsection{Responsabilidad en las simples asociaciones (art. 191)}

\begin{formula}{Responsabilidad en simples asociaciones (art. 191 CCyCN)}
Si la simple asociación no puede pagar a sus acreedores, los administradores y todos los miembros que hubieren participado de hecho en la administración responden solidariamente con sus propios bienes por las deudas de la simple asociación.

Los demás miembros que no participaron en la administración solo responden hasta el límite de las cuotas y contribuciones prometidas o debidas.
\end{formula}

En la simple asociación existe una menor diferencia entre la persona jurídica y sus miembros: quienes participaron de alguna forma en la administración y llevaron a un resultado de quebranto o insolvencia deben responder con sus propios bienes.

% ============================================================
\chapter{Cuadro Cornell de repaso general}
% ============================================================

El siguiente cuadro reúne, según el método Cornell, las preguntas centrales que atraviesan los cuatro bloques temáticos estudiados, con sus respuestas y una síntesis final para el repaso integral de la materia.

\begin{longtable}{
    >{\raggedright\arraybackslash}p{0.28\textwidth}
    >{\raggedright\arraybackslash}X
}
\cornellhead
\endfirsthead
\cornellhead
\endhead
\midrule
\multicolumn{2}{r}{\textit{continúa en la página siguiente}} \\
\midrule
\endfoot
\bottomrule
\endlastfoot

\textbf{¿Qué es una persona jurídica?} &
Son los entes a los cuales el ordenamiento jurídico les confiere aptitud para adquirir derechos y contraer obligaciones para el cumplimiento de su objeto y los fines de su creación (art. 141 CCyCN). A diferencia de la persona humana, cuya personalidad es reconocida por el legislador, la de las personas jurídicas es conferida por él. \\[0.4em]

\textbf{¿En qué consiste el principio de especialidad?} &
La capacidad de la persona jurídica está acotada a su objeto y sus fines específicos. No puede hacer cualquier cosa: solo aquello para lo que fue creada. De ahí la importancia de redactar el objeto con precisión y amplitud suficiente en el estatuto. \\[0.4em]

\textbf{¿Cuál es la diferencia entre la teoría de la ficción y las teorías de la realidad?} &
La teoría de la ficción (Savigny) entiende que la persona jurídica es un artificio del legislador, sin voluntad real. Las teorías de la realidad (organicista, de la institución) buscan un sustrato real. Las teorías normativistas las conciben como centros de imputación de normas. El CCyCN combina elementos normativistas e institucionales. \\[0.4em]

\textbf{¿Cómo clasifica el Código a las personas jurídicas?} &
Públicas: Estado nacional, provincias, municipios, entidades autárquicas, estados extranjeros, organismos internacionales, Iglesia Católica. Privadas: sociedades, asociaciones civiles, simples asociaciones, fundaciones, mutuales, cooperativas, consorcios de propiedad horizontal y otras previstas en leyes especiales. \\[0.4em]

\textbf{¿Qué norma prevalece según el artículo 150?} &
Orden de prevalencia: (1) normas imperativas de la ley especial; (2) normas imperativas del CCyCN; (3) estatuto y reglamentos; (4) normas supletorias. El estatuto tiene fuerza de norma jurídica en el ámbito de la persona jurídica. \\[0.4em]

\textbf{¿Cuándo se aplica la inoponibilidad del artículo 144?} &
Cuando la persona jurídica es usada para violar la ley, el orden público, la buena fe o para frustrar derechos de terceros, la personería se vuelve inoponible y se puede accionar directamente contra los socios, miembros o controlantes que hicieron posible esa actuación. Responden solidaria e ilimitadamente. \\[0.4em]

\textbf{¿Qué requisitos debe cumplir el nombre de una persona jurídica?} &
Debe indicar la forma jurídica adoptada; debe ser veraz, novedoso y con aptitud distintiva (diferenciarse de nombres, marcas y fantasías); no puede inducir a error sobre el objeto; no puede contener términos contrarios a la ley, el orden público o las buenas costumbres. Si incluye el nombre de una persona humana, se necesita su conformidad. \\[0.4em]

\textbf{¿Cuál es la diferencia entre domicilio y sede social?} &
Domicilio: ciudad indicada en el estatuto; determina jurisdicción y juez competente. Sede social: dirección física inscripta ante el organismo de contralor; las notificaciones allí enviadas se tienen por válidas. Domicilio especial: en el lugar de las sucursales, para las obligaciones allí contraídas. \\[0.4em]

\textbf{¿Cómo responde civilmente la persona jurídica?} &
Tres supuestos: (1) daños causados por directivos/administradores en ejercicio o con ocasión de sus funciones (art. 1763); (2) daños por cosas o actividades riesgosas, responsabilidad objetiva (art. 1757); (3) daños causados por dependientes en ejercicio o con ocasión de funciones (art. 1753). Los administradores también responden ilimitada y solidariamente ante la propia persona jurídica (art. 160). \\[0.4em]

\textbf{¿Cuáles son las causales de disolución del artículo 163?} &
Decisión de los miembros; condición resolutoria; consecución o imposibilidad del objeto; vencimiento del plazo; quiebra; fusión o escisión; reducción a uno del número de miembros; revocación o denegatoria de autorización estatal; agotamiento de bienes; otras causas estatutarias o legales. \\[0.4em]

\textbf{¿Qué garantías protegen frente a la revocación estatal del artículo 164?} &
Solo procede por actos graves que violen la ley, el estatuto o el reglamento. Requiere resolución fundada, garantía del derecho de defensa, y es apelable judicialmente con posible efecto suspensivo. El fundamento es la libertad de asociación constitucional. \\[0.4em]

\textbf{¿Qué caracteriza a las fundaciones?} &
Persona jurídica sin miembros, con fin de bien común y sin lucro, creada por el aporte patrimonial de uno o varios fundadores. Requiere autorización estatal para funcionar. Necesita patrimonio inicial suficiente (plan trienal). Órgano de gobierno: Consejo de Administración (mínimo tres personas). Solo fiscalización externa. El objeto es inmutable, salvo imposibilidad sobreviniente. \\[0.4em]

\textbf{¿Qué caracteriza a las asociaciones civiles?} &
Corporación de personas con fin común no lucrativo. Instrumento constitutivo: escritura pública. Tres órganos: Comisión Directiva, Asamblea (suprema) y órgano de fiscalización (Comisión Revisora de Cuentas si hay más de 100 socios). Poder disciplinario: proceso con derecho de defensa, gradualidad y proporcionalidad. Apelación judicial solo sobre la forma, salvo grave violación de derechos humanos. \\[0.4em]

\textbf{¿En qué se diferencian las simples asociaciones?} &
No requieren autorización estatal. Pueden constituirse por instrumento privado con firma certificada. Comienzan su existencia desde la constitución. Si hay menos de 20 socios, pueden prescindir del órgano de fiscalización. Principal diferencia: en caso de insolvencia, los administradores y miembros que participaron en la gestión responden solidariamente con sus propios bienes. \\[0.4em]

\midrule
\multicolumn{2}{p{0.94\textwidth}}{
\textbf{Resumen:}
Las personas jurídicas son entes colectivos a los que el ordenamiento jurídico --no la naturaleza-- otorga aptitud para ser titulares de derechos y obligaciones, con el fin de cumplir el objeto para el que fueron creadas (principio de especialidad). El Código Civil y Comercial las clasifica en públicas --el Estado y sus organismos, los estados extranjeros, los organismos internacionales y la Iglesia Católica-- y privadas --sociedades, asociaciones civiles, simples asociaciones, fundaciones, mutuales, cooperativas y consorcios--. Su existencia comienza desde la constitución, aunque algunas requieren autorización estatal para funcionar. Son entidades con personalidad diferenciada de sus miembros (patrimonios y responsabilidades distintas), salvo los supuestos de inoponibilidad del artículo 144, donde la personería cede frente al uso abusivo o fraudulento. Responden civilmente por los daños que causen sus directivos en ejercicio de sus funciones (art. 1763), por las cosas y actividades riesgosas que posean (art. 1757) y por los hechos dañosos de sus dependientes (art. 1753); además, los propios administradores responden ilimitada y solidariamente frente a la persona jurídica (art. 160). Su disolución puede deberse a la voluntad de los miembros, al cumplimiento o imposibilidad del objeto, al vencimiento del plazo, a la quiebra o a la revocación estatal --esta última con estrictas garantías de fundamentación y defensa--. Entre los tipos especiales, las fundaciones se distinguen por la ausencia de miembros, la afectación de un patrimonio a un fin altruista y el control externo exclusivo de la autoridad de contralor; las asociaciones civiles, por su carácter corporativo, sus tres órganos de gobierno y el poder disciplinario sobre sus miembros con plenas garantías procesales; y las simples asociaciones, por no requerir autorización estatal y por la responsabilidad solidaria que recae sobre los administradores y miembros activos ante la insolvencia.
} \\

\end{longtable}

\end{document}
